\documentclass[10pt,a4paper]{amsart}
\usepackage[margin=19mm]{geometry}
\usepackage{amsmath,amssymb,mathtools}
\usepackage[scr=rsfs,cal=euler]{mathalfa}
\usepackage{xfrac}
\usepackage[unicode,hidelinks,bookmarks=false]{hyperref}
\newcommand{\ie}{i.e.\ }
\newcommand{\cf}{cf.\ }
\newcommand{\resp}{respectively\ }
\newcommand{\Subsection}{\subsection}
\providecommand{\nobreakdash}{\nobreak\hbox{-}}
\setlength{\emergencystretch}{2em}
\multlinegap0pt
\usepackage{bm}
\usepackage{bbm}
\usepackage{dsfont}
\usepackage{xspace}
\usepackage{cancel}
\usepackage{mathtools}
\usepackage{amsthm}
\makeatletter
\@ifundefined{theorem}{\newtheorem{theorem}{Theorem}}{}
\@ifundefined{corollary}{\newtheorem{corollary}[theorem]{Corollary}}{}
\@ifundefined{lemma}{\newtheorem{lemma}[theorem]{Lemma}}{}
\makeatother
\DeclareMathOperator {\mono} {mp}
\DeclareMathOperator{\cp}{\,\square\,}
\title[Monophonic position sets of Cartesian and lexicographic prod]{Monophonic position sets of Cartesian and lexicographic products of graphs}
\author{Ullas Chandran S. V., Sandi Klav\v{z}ar, Neethu P. K., James Tuite}
\date{Source: arXiv:2412.09837v3 (CC BY 4.0)}
\begin{document}
\maketitle

\noindent\emph{Source and licence.} Monophonic position sets of Cartesian and lexicographic products of graphs. Version arXiv:2412.09837v3 (CC BY 4.0), distributed under the Creative Commons Attribution 4.0 International licence, as recorded in the arXiv licence metadata for this exact version and shown on the abstract page https://arxiv.org/abs/2412.09837v3. Full rights notice: licenses/arxiv-2412.09837-CC-BY-4.0.txt.\par\smallskip

\noindent\textit{Editorial scope.} Counted unit: the Theorem whose source label is \texttt{thm:bipartite} in the article (source0.tex lines 631--635, source label \texttt{thm:bipartite}), delivered together with the article's own complete original proof of it (source0.tex lines 637--651, one \texttt{proof} environment of 15 lines). No mathematics is written, repaired or re-derived: statement and proof are the article's own text line for line. Required context retained uncounted: cor:varied sets are small (lines 477--479); lem:simplicial vertices (lines 589--592); thm:general bounds (lines 604--608). Every definition, display and label of those lines is retained unchanged. Omitted from this package and not redistributed: 1--476, 480--588, 593--603, 609--630, 636--636, 652--981. Nothing inside a retained excerpt is omitted or altered: every retained body is delivered exactly as the article's lines are written, so no omission receipt of any kind accompanies this package. Citations: the retained excerpts carry no citation command at all, so no bibliography entry of the article is transcribed and no reference is redistributed. Reference adaptation: none was needed and none was made; every reference inside the delivered unit resolves inside it, so no unresolved reference or citation is produced. Printed slips kept literal: no printed word, symbol, hypothesis or number of the counted statement or of its proof was altered. Layout and class: the article's own class is \texttt{article} with packages including amsmath, amsthm, amscd, amsfonts, amssymb, graphicx, latexsym, mathtools, fontenc, tikz, color, url, todonotes, cite; the wrapper is the lane's pinned \texttt{amsart} editorial wrapper at 10pt on A4, which restates the theorem-like environments the retained text needs and adds the article's own macro definitions that the retained text needs (\texttt{\textbackslash cp}, \texttt{\textbackslash mono}), with extra wrapper packages (bm, bbm, dsfont, xspace, cancel, mathtools). The delivered numbering is the unit's own. Assets: no retained span contains a figure environment, a graphics-inclusion command, a PDF-page inclusion, a table or a TikZ picture; no image file is copied into this package, the package declares no asset dependency and no image file is redistributed with it. Licence: version arXiv:2412.09837v3 of this article is distributed under the Creative Commons Attribution 4.0 International licence, as recorded in the arXiv licence metadata for this exact version and shown on the abstract page https://arxiv.org/abs/2412.09837v3. A full rights notice is delivered at licenses/arxiv-2412.09837-CC-BY-4.0.txt. Characters: every retained excerpt is pure ASCII, so the delivered engine, XeLaTeX with its default Latin Modern text font, reports no missing character.
\begin{corollary}\label{cor:varied sets are small}
If $S$ is a varied monophonic position set of $G \cp H$, then $|S| \leq 2$.
\end{corollary}

\begin{lemma}
\label{lem:simplicial vertices}
If $S$ is a layered or cliquey monophonic position set of $G \cp H$ (where $\pi _G(S)$ is a clique) and $|S| > \max \{ \omega (G),\omega (H)\} $, then every vertex from $\pi _G(S)$ is simplicial.
\end{lemma}

\begin{theorem} 
\label{thm:general bounds}
If $G$ and $H$ are connected graphs, then $$\max\{\omega(G), \omega(H)\}\leq \mono (G \cp H) \leq \max \{ \omega (G), \omega (H), \sigma (G) \mono_i(H), \sigma (H) \mono_i(G)\}.$$ 
    Furthermore, if neither $G$ nor $H$ has simplicial vertices, then $$\mono (G \cp H) = \max \{ \omega (G), \omega (H)\}\,.$$    
\end{theorem}

\begin{theorem}
\label{thm:bipartite}
If $G$ and $H$ are connected triangle-free graphs of order at least $3$, then 
$$\mono(G\cp H)\leq\max\{2, \sigma(G)\Delta(H), \sigma(H)\Delta(G)\}\,.$$ 
\end{theorem}

\begin{proof} 
Let $S$ be a largest monophonic position set of $G\cp H$. If $|S|\geq 3$, then $S$ is either cliquey or layered by Corollary \ref{cor:varied sets are small}. Suppose that $S$ is cliquey and $\pi_G(S)=\{u_1, u_2\}$ induces a clique, that is, $u_1\sim u_2$ in $G$. However, as the order of $G$ is at least three, the vertices $u_1$ and $u_2$ cannot both be leaves, contradicting Lemma~\ref{lem:simplicial vertices}. 

Hence $S$ must be layered. Again by Lemma~\ref{lem:simplicial vertices}, we may assume that $\pi_G(S)=\{u\}$, where $u$ is a leaf of $G$. Let $u^{\prime }$ be any vertex of $G$ with $d_G(u,u^{\prime }) = 2$ and $P$ be any shortest $u,u^{\prime }$ path. As in the proof of Theorem~\ref{thm:general bounds}, we conclude that $\pi_H(S)$ is an independent monophonic position set of $H$. 

We first show that all vertices of $\pi _H(S)$ are at distance two from each other. Let $k = \max \{ d_H(v,v^{\prime }):v,v^{\prime } \in \pi _H(S)$. Assume for a contradiction that $k \geq 3$ and let $v_1,v_2$ be a pair of vertices of $\pi _H(S)$ with $d_H(v_1,v_2) = k$. Let $v_3$ be any third vertex of $\pi _H(S)$ ($v_3$ exists, since $|\pi _H(S)| = |S| \geq 3$). Amongst all shortest $v_1,v_2$-paths in $H$, let $Q$ be a path that minimises $d_H(Q,v_3)$, where $d_H(Q,v_3) = \min \{ d_H(x,v_3):x \in V(Q)\} $. Let $Q$ be the path $v_1 = x_0,x_1,\dots x_{k-1},x_k = v_2$. 

Let $x_i$ be the vertex with the smallest value of $i$ on $Q$ such that $d_H(x_i,v_3) = d_H(Q,v_3)$. Let $Q^+$ be the monophonic $v_3,v_1$-path in $H$ formed from the concatenation of a shortest $v_3,x_i$-path $Q^*$ and the $x_i,v_1$-section of $Q$. If $i = k$, then $Q^+$ would be a monophonic path containing all three of $v_1,v_2,v_3$, so we have $0 \leq i \leq k-1$.   

Suppose that $0 \leq i \leq k-2$. Consider the concatenation of the four paths $_uQ^+$, $P_{v_1}$, $_{u^{\prime }}Q$ and $\widetilde{P}_{v_2}$ in that order. As $v_2$ does not lie on $Q^+$, $d_G(u,u^{\prime }) = 2$ and $v_1 \not \sim v_2$ in $H$, this concatenated path would be a monophonic path containing $(u,v_1)$, $(u,v_2)$ and $(u,v_3)$ unless $v_2$ has a neighbour $y$ on $Q^*$. By our assumption that $i \leq k-2$, $y \neq x_i$. In fact, $y$ must be the neighbour of $x_i$ on $Q^*$, for otherwise $Q^*$ would not be a shortest path from $v_3$ to $Q$. As $Q$ is a shortest $v_1,v_2$-path, it follows that $i = k-2$ and the path formed by following the $v_1,v_{k-2}$-section of $Q$ followed by the path $v_{k-2},y,v_2$ would be a shortest $v_1,v_2$-path closer to $v_3$ than $Q$, contradicting the definition of $Q$. Hence $i = k-1$.   

Now the vertex $v_2$ has no neighbours on $Q^*$ other than $x_{k-1}$, for otherwise $H$ would either contain a triangle or else $Q^*$ would not be a shortest path from $v_3$ to $Q$. Therefore the $v_3,v_2$-path $Q^{++}$ formed by appending the edge $x_{k-1}v_2$ to $Q^*$ is monophonic. By definition of $Q^*$ and the condition that $d_H(v_1,v_2) = k \geq 3$ there is no edge from $v_1$ to $Q^{++}$ and it follows that the concatenation of the four paths $_uQ^{++}$, $P_{v_2}$, $_{u^{\prime }}\widetilde{Q}$ and $\widetilde{P}_{v_1}$ is a monophonic path containing each of $(u,v_1)$, $(u,v_2)$ and $(u,v_3)$. Therefore our assumption that $k \geq 3$ is false and all pairs of vertices from $\pi _H(S)$ are at distance two.

Suppose now that $\mono (G \cp H) > \Delta (H)$. Fix two vertices $v_1,v_2 \in \pi _H(S)$ and let $Q$ be a shortest path $v_1,x,v_2$ between them. As $|\pi _H(S)| > \Delta (H)$, there must be a vertex $v_3$ of $\pi _H(S)$ that is not adjacent to $x$. Let $Q'$ be a shortest path $v_2,y,v_3$. As there is no edge from $v_3$ to $\{ v_1,x,v_2\} $, it follows that the concatenation of $_uQ$, $P_{v_2}$, $_{u^{\prime }}Q'$ and $\widetilde{P}_{v_3}$ in that order would be a monophonic path containing $(u,v_1)$, $(u,v_2)$ and $(u,v_3)$, a contradiction. Therefore $\mono (G \cp H) \leq \Delta (H)$. The claimed inequality follows when we allow the roles of $G$ and $H$ to be reversed.
\end{proof}

\end{document}
